\chapter{Conclusión}

En este trabajo se utilizaron los métodos de \acrshort{rl}: \acrshort{gps} (\textit{model-based}) y \acrshort{ddpg} (\textit{model-free}) para optimizar diferentes \acrshort{ann} con el propósito de controlar un cuadricoptero, llevarlo al origen alineado con los ejes, con velocidades angulares y lineales iguales a cero, tal como lo hace un controlador \acrshort{ilqr}. De hecho, se uso el método \acrshort{ilqr} para conseguir una configuración subóptima del costo, que fue fundamental para los métodos de \acrshort{rl}. El primer entrenamiento se realizó con el método \acrshort{gps}, debido a que el uso de controladores locales \acrshort{ilqr} generaba mayor certeza en los resultados. Esto a su vez permitió definir una arquitectura de la \acrshort{ann} asociada a la política e hiperparametros que después serian retomados en el método \acrshort{ddpg}. 

Para el método \acrshort{ddpg} se utilizó una configuración mucho más flexible que el costo de \acrshort{gps} con la finalidad de recompensar un comportamiento deseado y penalizar uno no deseado. 

Una vez entrenadas las diferentes políticas con los métodos \acrshort{gps} y \acrshort{ddpg} se definieron conjuntos para probar la estabilidad y distintos al usado en el entrenamiento para determinar la confiabilidad de estos métodos. Esto dio lugar a una comparación no estricta de los métodos basados en redes neuronales con los métodos de control más convencionales: \acrshort{ilqr} y control linealizado por retroalimentación. Se observa en las figuras \ref{fig:stability_region_ilqr} y \ref{fig:stability_region_linear} que bajo sus respectivos criterios ambos controladores convencionales generan regiones conexas de posiciones iniciales cuyas trayectorias correspondientes terminaron cerca del objetivo. 

La política resultante de \acrshort{gps} imitó el comportamiento del controlador \acrshort{ilqr} todas las variables salvo en el ángulo $\psi$ y velocidad $r$ en muestras de la distribución $\mathcal{X}$, como se observa en \ref{fig:state_traj_gps}. Al excluir esas variables, la política genera una cantidad de trayectorias estables, incluso varias fuera de las regiones conexas de los métodos convencionales. Se destaca que en el eje $(r-\psi)$ se observa un alto porcentaje de estabilidad en los rangos de perturbación definidos, mientras que en los ejes angulares $(p-\varphi)$ y $(q-\theta)$ son mínimas las trayectorias estables. Sin embargo, en general no se encontraron regiones conexas que permitan determinar la confiabilidad de la política. 

El método \acrshort{ddpg} falla en la tarea fijada en este trabajo, esto puede deberse al mecanismo descrito en \textit{The problem with DDPG: understanding failures in deterministic environments with sparse rewards} \cite{Matheron2020}. El algoritmo parece entrar en punto fijo subóptimo. La región de acciones que estabilizan el cuadricoptero es muy estrecha y la recompensa, aunque formalmente es densa, es efectivamente escasa (debido a $\tanh^2$). Esto fomenta que el actor converja rápidamente a una política de acciones constantes antes de haber observado transiciones relevantes. Entonces el crítico converge a una función de valor plana. El sistema entra así en un \textit{deadlock} actor-crítico descrito en \cite{Matheron2020}. El actor no puede cambiar porque el gradiente es nulo, y el crítico no generaliza a mejores acciones, esto hace que \acrshort{ddpg} quede atrapado en un política incapaz de estabilizar el cuadricoptero. Esta caracterización se hace patente en las acciones mostradas en la figura \ref{fig:ddpg_actions}. Se observa que con $10$ \textit{epochs} el entrenamiento no es suficiente para obtener una política efectiva; sin embargo, con una mayor cantidad de epochs–como 200– la política sigue siendo inefectiva y es prácticamente constante. 

Este trabajo echibe que los métodos \acrshort{ddpg} y \acrshort{gps} son flexibles, fuertemente adaptativos a las nuevas experiencias, potencialmente útiles para aprender tareas complejas como se demuestra \cite{hwangbo2017control}, \cite{zhang2016}. Precisamente estos dos trabajos persiguen un enfoque híbrido entre \acrshort{rl}, el uso de redes neuronales y métodos más convencionales de control. Por si sólo \acrshort{gps} ya es un enfoque híbrido. Sin embargo, en el caso del control de un cuadricoptero, las políticas entrenadas no es suficiente para generar trayectorias estables en diferentes perturbaciones, esto puede deberse a la la complejidad de la dinámica. A pesar de ello este trabajo sirve preámbulo para adentrarse comprender el método \acrshort{mpc}-\acrshort{gps} \cite{zhang2016}, que resuelve el problema de generar trayectorias y evitar objetos por parte de un cuadricoptero. Esto sugiere que este método podría llevar acabo con mejor efectividad las pruebas de estabilidad presentadas aquí. También este trabajo muestra las limitaciones del método \acrshort{ddpg} para entrenar por si sólo una política que cumpla el objetivo planteado. Sin embargo, esto puede dar lugar a usar métodos más robustos y estables como \acrshort{td3} \cite{fujimoto2018} o \acrshort{sac} \cite{haarnoja2018}, así como explorar estrategias híbridas preservando el enfoque \textit{model-free}. 

\appendix

